\section{Foundations of Quantum Mechanics}

We start with an overview of corresponding ideas between linear algebra and quantum mechanics.

\begin{table}[ht!]
    \centering
    \begin{tabular}{c|c}
        Linear Algebra & Quantum Mechanics \\
        \hline
        Vector (\(n\)-dimensional complex vector space) & State (\(\infty\)-dimensional vector space) \\
        \(\vb{v}, \vb{w}, \ldots\) & \(\psi, \phi, \ldots\) \\
        Discrete Basis \(\{\vb{e}_i\}_{i=1}^n\) & Continuous Basis \(\Bqty{\vb{x}}_{\vb{x} \in \RR^3}\) \\
        Described using \(\vb{v} \mapsto \pmqty{v_1 \\ \vdots \\ v_n}\) & Described using \(\vb{x} \mapsto \psi\pqty{\vb{x}, t}\) \\
        Denoted \(\CC^n\) & Denoted \(L^2\pqty{\RR^3}\) \\
        Inner Product \(\pqty{\vb{v}, \vb{w}} = \sum_{i=1}^n {v_i\conjstar} w_i\) & Inner Product \(\pqty{\psi, \phi} = \int_{\RR^3} \dd[3]{\vb{x}} \psi\conjstar \pqty{\vb{x}, t} \phi\pqty{\vb{x}, t}\) \\
        Linear Map \(\functiondomain{T}{\CC^n}{\CC^n}\) & Operator \(\functiondomain{\hat{O}}{L^2\pqty{\RR^3}}{L^2\pqty{\RR^3}}\) \\
        Matrix Representation \(\pmqty{T_{ij}}_{i,j=1}^n\) & Operator \(\hat{O}\) \\
        \(\vb{w} = T \vb{v}\) & \(\phi = \hat{O} \psi\) \\
    \end{tabular}

\end{table}

\subsection{State, Wavefunction and Probability}

In quantum mechanics, the wavefunction \(\psi\) determines the dynamics of particles in a probabilistic way.

\begin{definition}[State, Wavefunction]
    \(\psi\) is the \term{state} of a particle, and \(\psi \pqty{\vb{x}, t}\) is the complex coefficient of \(\psi\) in the \emph{continuous} basis \(\vb{x}\) at a given time \(t\), i.e., \(\psi\pqty{\vb{x}, t}\) is \(\psi\) in the \(\vb{x}\)-representation, and is called the \term{wavefunction}:
    \begin{equation}
        \functiondomain{\psi\pqty{\vb{x}, t}}{\RR^3}{\CC} \quad \forall t \in \RR.
    \end{equation}
\end{definition}

\term{Born's rule} of \term{probabilistic interpretation} gives a physical interpretation of the wavefunction, hence leading to some mathematical property that we will see. The \term{probability density} for a particle described by \(\psi\) to sit at \(\vb{x}\) at a given time \(t\) is given by
\begin{equation}
    \rho\pqty{\vb{x}, t} \propto \abs{\psi\pqty{\vb{x}, t}}^2 = \psi\conjstar \pqty{\vb{x}, t} \psi\pqty{\vb{x}, t}.
\end{equation}
and the probability that the particle lies in some infinitesimal volume \(\dd{V}\) around \(\vb{x}\) is given by
\begin{equation}
    \rho\pqty{\vb{x}, t} \dd{V}.
\end{equation}

This gives us the dimension of the wavefunction \(\bqty{\psi} = \dimL^{-\flatfrac{3}{2}}\). In addition, this probabilistic interpretation leads to the following mathematical properties.

\begin{enumerate}[label=(\roman*)]
    \item \term{Square-integrability}/\term{normalisable}: a wavefunction \(\psi\) must be square-integrable, i.e. for some constant \(N \in \RR\),
    \begin{equation}
        \int_{\RR^3} \dd[3]{\vb{x}} \abs{\psi\pqty{\vb{x}, t}}^2 = N < \infty \quad \forall t \in \RR.
    \end{equation}

    \item The wavefunction must be \term{normalised}, i.e. \(N = 1\), so that the total probability of finding the particle somewhere in space is \(1\): if some wavefunction \(\psi\) is not normalised, we can always normalise it by defining
    \begin{equation}
        \overline{\psi}\pqty{\vb{x}, t} = \frac{1}{\sqrt{N}} \psi\pqty{\vb{x}, t}.
    \end{equation}
    then
    \begin{equation}
        \int_{\RR^3} \dd[3]{\vb{x}} \abs{\overline{\psi}\pqty{\vb{x}, t}}^2 = 1 \quad \forall t \in \RR.
    \end{equation}

    When \(\psi\) is normalised, we have
    \begin{equation}
        \rho\pqty{\vb{x}, t} = \abs{\psi\pqty{\vb{x}, t}}^2.
    \end{equation}

    It is common for us to denote \(\psi\) for both \(\psi\) and \(\overline{\psi}\), the normalised version.

    \item If the wavefunction \(\tilde{\psi} \pqty{\vb{x}, t} = \exp\pqty{i \alpha} \psi\pqty{\vb{x}, t}\) for some constant \(\alpha \in \RR\), then notice that the probabilities are the same, i.e. \(\abs{\tilde{\psi}\pqty{\vb{x}, t}}^2 = \abs{\psi\pqty{\vb{x}, t}}^2\). It turns out that this does not affect anything else, and so \(\tilde{\psi}\) is equivalent to \(\psi\).
\end{enumerate}

\begin{remark}
    If \(\alpha = \alpha\pqty{\vb{x}}\) is space-dependent, then \(\tilde{\psi}\) is not equivalent to \(\psi\): although the two give the same probabilities, this changes other observables, and hence are not equivalent.
\end{remark}

In other words, states in quantum mechanics correspond to equivalent classes of wavefunctions, where two wavefunctions are equivalent if they differ by
\[
    \tilde{\psi} \pqty{\vb{x}, t} = \lambda \psi\pqty{\vb{x}, t}
\]
for any \(\lambda \in \CC\) with \(\lambda \neq 0\), combining the normalisation condition and the phase ambiguity.

\subsection{The Hilbert Space \(L^2\pqty{\RR^3}\)}

\begin{definition}[Hilbert Space \(\hilbertSpace\)]
    The set of all square-integrable functions \(\psi\) on \(\RR^3\) forms a Hilbert space, denoted \(\hilbertSpace = L^2\pqty{\RR^3}\).
\end{definition}

\begin{theorem}[Square-Integrability of Linear Combination]
    If \(\psi_1\) and \(\psi_2\) are both square-integrable functions, then their linear combination \(\psi = a_1 \psi_1 + a_2 \psi_2\) is also square-integrable for any constants \(a_1, a_2 \in \CC\).
\end{theorem}

\begin{proof}
    Let \(N_1, N_2 \in \RR\) be such that
    \begin{equation}
        \int_{\RR^3} \dd[3]{\vb{x}} \abs{\psi_1\pqty{\vb{x}, t}}^2 = N_1 < \infty, \quad
        \int_{\RR^3} \dd[3]{\vb{x}} \abs{\psi_2\pqty{\vb{x}, t}}^2 = N_2 < \infty.
    \end{equation}

    Hence, we have
    \begin{subequations}
        \begin{align}
            \int_{\RR^3} \dd[3]{\vb{x}} \abs{\psi\pqty{\vb{x}, t}}^2 &= \int_{\RR^3} \dd[3]{\vb{x}} \abs{a_1 \psi_1 + a_2 \psi_2}^2 \\
            &\leq \int_{\RR^3} \dd[3]{\vb{x}} \pqty{\abs{a_1 \psi_1} + \abs{a_2 \psi_2}}^2 \\
            &= \int_{\RR^3} \dd[3]{\vb{x}} \pqty{\abs{a_1 \psi_1}^2 + \abs{a_2 \psi_2}^2 + 2 \abs{a_1 \psi_1} \abs{a_2 \psi_2}} \\
            &\leq \int_{\RR^3} \dd[3]{\vb{x}} \pqty{\abs{a_1 \psi_1}^2 + \abs{a_2 \psi_2}^2 + \abs{a_1 \psi_1}^2 + \abs{a_2 \psi_2}^2} \\
            &= \int_{\RR^3} \dd[3]{\vb{x}} \pqty{2 \abs{a_1 \psi_1}^2 + 2 \abs{a_2 \psi_2}^2} \\
            &= 2 \abs{a_1}^2 N_1 + 2 \abs{a_2}^2 N_2 \\
            &< \infty.
        \end{align}
    \end{subequations}
    as desired.
\end{proof}

\subsection{Inner Product on \(\hilbertSpace\)}

\begin{definition}[Inner Product on \(\hilbertSpace\)]
    The \term{inner product} on \(\hilbertSpace\) is defined as
    \begin{equation}
        \pqty{\psi, \phi} = \int_{\RR^3} \dd[3]{\vb{x}} \psi\conjstar \pqty{\vb{x}, t} \phi\pqty{\vb{x}, t}.
    \end{equation}
    with \(\psi, \phi \in \hilbertSpace\).
\end{definition}

\begin{theorem}[Convergence of Inner Product]
    If \(\psi, \phi \in \hilbertSpace\), then the inner product \(\pqty{\psi, \phi}\) exists.
\end{theorem}

\begin{proof}
    Suppose we have
    \begin{subequations}
        \begin{align}
            \int_{\RR^3} \dd[3]{\vb{x}} \abs{\psi\pqty{\vb{x}, t}}^2 &= N_1 < \infty, \\
            \int_{\RR^3} \dd[3]{\vb{x}} \abs{\phi\pqty{\vb{x}, t}}^2 &= N_2 < \infty,
        \end{align}
    \end{subequations}
    then
    \begin{subequations}
        \begin{align}
            \abs{\pqty{\psi, \phi}} &= \abs{\int_{\RR^3} \dd[3]{\vb{x}} \psi\conjstar \phi} \\
            &\leq \sqrt{\int_{\RR^3} \dd[3]{\vb{x}} \abs{\psi\pqty{\vb{x}, t}}^2} \sqrt{\int_{\RR^3} \dd[3]{\vb{x}} \abs{\phi\pqty{\vb{x}, t}}^2} \\
            &= \sqrt{N_1} \sqrt{N_2} < \infty,
        \end{align}
    \end{subequations}
    by the integral Schwarz inequality, as desired.
\end{proof}

\begin{proposition}[Properties of The Inner Product]
    \begin{enumerate}
        \item \emph{Hermitian}: \(\pqty{\psi, \phi} = \pqty{\phi, \psi}\conjstar\), for all \(\psi, \phi \in \hilbertSpace\).
        \item \emph{Antilinear in First Entry, Linear in Second Entry}:
        \begin{subequations}
            \begin{align}
                \pqty{a_1 \psi_1 + a_2 \psi_2, \phi} &= a_1\conjstar \pqty{\psi_1, \phi} + a_2\conjstar \pqty{\psi_2, \phi}, \\
                \pqty{\psi, b_1 \phi_1 + b_2 \phi_2} &= b_1 \pqty{\psi, \phi_1} + b_2 \pqty{\psi, \phi_2},
            \end{align}
        \end{subequations}
        for all \(a_1, a_2 \in \CC, \psi, \psi_1, \psi_2, \phi, \phi_1, \phi_2 \in \hilbertSpace\).
        \item \emph{Inner product of a function with itself is positive}:
        \begin{equation}
            \pqty{\psi, \psi} = \int_{\RR^3} \dd[3]{\vb{x}} \abs{\psi\pqty{\vb{x}, t}}^2 > 0,
        \end{equation}
        since the wavefunction \(\psi \equiv 0\) is not a valid state.
    \end{enumerate}
\end{proposition}

\begin{definition}[Norm of a Waverunction]
    The \term{norm} of a wavefunction \(\psi\) is defined as
    \begin{equation}
        \norm{\psi} = \sqrt{\pqty{\psi, \psi}} = \sqrt{\int_{\RR^3} \dd[3]{\vb{x}} \abs{\psi\pqty{\vb{x}, t}}^2}.
    \end{equation}
\end{definition}

\begin{definition}[Normalised Wavefunction]
    A wavefunction \(\psi\) is \term{normalised} if \(\norm{\psi} = 1\).
\end{definition}

\begin{definition}[Orthogonal Wavefunctions]
    Two wavefunctions \(\psi\) and \(\phi\) are \term{orthogonal} if \(\pqty{\psi, \phi} = 0\).
\end{definition}

\begin{definition}[Orthonormal Set of Wavefunctions]
    A set of wavefunctions \(\Bqty{\psi_n}\) is \term{orthonormal} if
    \begin{equation}
        \pqty{\psi_m, \psi_n} = \kronecker_{mn} = \begin{cases}
            1, & m = n, \\
            0, & m \neq n.
        \end{cases}
    \end{equation}
\end{definition}

\begin{definition}[Complete Set of Wavefunctions]
    A set of wavefunctions \(\Bqty{\psi_n}\) is \term{complete} if any wavefunction \(\psi\) can be expressed as a linear combination of the set, i.e., for all \(\psi \in \hilbertSpace\), there exists a set of constants \(c_n \in \CC\) such that
    \begin{equation}
        \psi = \sum_n c_n \psi_n.
    \end{equation}
\end{definition}

\begin{lemma}[Coefficients of Linear Combination]
    If a set of wavefunctions \(\Bqty{\psi_n}\) forms a complete and orthonormal basis of \(\hilbertSpace\), then the coefficients are given by
    \begin{equation}
        c_n = \pqty{\psi_n, \psi}.
    \end{equation}
\end{lemma}

\begin{proof}
    We have
    \begin{subequations}
        \begin{align}
            \pqty{\psi_n, \phi} &= \pqty{\psi_n, \sum_{m = 0}^{\infty} c_m \psi_m} \\
            &= \sum_{m = 0}^{\infty} c_m \pqty{\psi_n, \psi_m} \\
            &= \sum_{m = 0}^{\infty} c_m \kronecker_{nm} \\
            &= c_n.
        \end{align}
    \end{subequations}
\end{proof}

\subsection{Time-Dependent Schr\"odinger Equation}

The second postulate in quantum mechanics is the \term{time-dependent Schr\"odinger equation}, which reads
\begin{equation}
    \label{eq:time-dependent-schrodinger-equation}%
    i\hbar \pdv{t} \psi(\vb{x}, t) = \hat{H} \psi(\vb{x}, t),
\end{equation}
where the \term{Hamiltonian operator} \(\hat{H}\) is given by
\begin{equation}
    \label{eq:hamiltonian-operator}%
    \hat{H} = -\frac{\hbar^2}{2m} \laplacian + U(\vb{x}),
\end{equation}
where \(U \in \RR\) is the \term{potential energy} of the particle.

Note this equation has a first-derivative in \(t\) while a second derivative in \(\vb{x}\). This is a non-relativistic equation (but it has a relativistic generalisation).

Once the wavefunction \(\psi\pqty{\vb{x}, t_0}\) is known at a given time \(t_0\) at all points \(\vb{x}\), then the wavefunction is completely determined at all times \(t > t_0\).

We attempt to solve the Schr\"odinger equation for an \(\Pelectron\). Recall a general solution to the wave equation is given by \(\exp\pqty{i\pqty{kx - \omega t}}\), where \(\omega = ck\). Plugging this into \autoref{eq:time-dependent-schrodinger-equation} with \(U = 0\), we have
\begin{equation}
    - i \hbar \pdv{t} \psi\pqty{x, t} = - \frac{\hbar^2}{2m} \pdv[2]{x} \psi\pqty{x, t},
\end{equation}
and we get \(\omega = \frac{\hbar k^2}{2m}\), i.e.,
\begin{equation}
    \omega \propto k^2.
\end{equation}

This is consistent with the de Broglie waves in \autoref{eq:de-broglie-waves}: since for a free particle, \(E = \flatfrac{p^2}{2m}\), we have
\begin{equation}
    \omega = \frac{E}{\hbar} = \frac{p^2}{2 m \hbar} = \frac{\hbar k^2}{2 m}.
\end{equation}

\begin{proposition}[Norm is Constant in Time]
    If \(\psi\) obeys \autoref{eq:time-dependent-schrodinger-equation}, then the norm of \(\psi\) is constant in time.
\end{proposition}

\begin{proof}
    Let the norm be
    \begin{equation}
        N(t) = \int_{\RR^3} \dd[3]{\vb{x}} \abs{\psi\pqty{\vb{x}, t}}^2.
    \end{equation}

    We want to show that \(\dot{N} = 0\). Notice
    \begin{subequations}
        \begin{align}
            \dv{N}{t} &= \dv{t} \int_{\RR^3} \dd[3]{\vb{x}} \abs{\psi\pqty{\vb{x}, t}}^2 \\
            &= \int_{\RR^3} \dd[3]{\vb{x}} \pdv{t} \abs{\psi\pqty{\vb{x}, t}}^2.
        \end{align}
    \end{subequations}

    Note that we differentiate to get
    \begin{equation}
        \pdv{t} \pqty{\psi\conjstar \psi} = \psi\conjstar \pdv{\psi}{t} + \psi \pdv{\psi\conjstar}{t}.
    \end{equation}

    By the Schr\"odinger equation, we have
    \begin{subequations}
        \begin{align}
            \pdv{\psi}{t} &=\frac{i \hbar}{2m} \laplacian \psi - \frac{i}{\hbar} U \psi, \\
            \pdv{\psi\conjstar}{t} &= -\frac{i \hbar}{2m} \laplacian \psi\conjstar + \frac{i}{\hbar} U \psi\conjstar.
        \end{align}
    \end{subequations}

    Then, note that
    \begin{equation}
        \pdv{t} \pqty{\psi\conjstar \psi} = \div \pqty{\frac{i \hbar}{2m} \pqty{\psi\conjstar \grad{\psi} - \psi \grad{\psi\conjstar}}},
    \end{equation}
    and hence
    \begin{equation}
        \frac{N}{t} = \int_{\RR^3} \dd[3]{\vb{x}} \div \bqty{\frac{i \hbar}{2m} \pqty{\psi\conjstar \grad{\psi} - \psi \grad{\psi\conjstar}}} = 0
    \end{equation}
    by the divergence theorem, since \(\psi\) is square-integrable and hence vanishes at infinity.
\end{proof}